\chapter{Integración Terraform de la Fase 2}

La infraestructura nueva se separa en archivos con responsabilidades claras:

\begin{itemize}
  \item \code{dmz.tf}: tabla y asociación de ruta de la DMZ.
  \item \code{nacls.tf}: NACL y asociaciones de Visitas, TI y DMZ.
  \item \code{phase2-security.tf}: Security Groups base de servicios y pruebas.
  \item \code{locals.tf}: DMZ, reserva, direcciones y etiquetas.
  \item \code{variables.tf}: dominios, IP, HTTPS y control de cómputo.
  \item \code{outputs.tf}: identificadores y datos de integración.
\end{itemize}

La variable \code{enable_phase2_compute} utiliza \code{false} como valor predeterminado.
Agregarla no habilita recursos por sí mismo: los archivos de cómputo de los demás responsables
deben condicionar sus instancias con esta variable.

La validación local se ejecutó con Terraform 1.17.2 y el proveedor oficial AWS 6.73.0. El
formato recursivo y \code{terraform validate} finalizaron correctamente. Además,
\path{scripts/validate-pablo-phase2.py} verificó contención, no solapamiento, capacidad,
direcciones asignables y presencia de los controles estructurales esperados.

\section{Plan y despliegue observado en AWS}

Durante la ventana autorizada del 17 de noviembre de 2026 se generó un plan guardado con
\textbf{41 recursos por agregar, 0 por modificar y 0 por destruir}. La revisión mostró las siete
redes del diseño, los clientes temporales y las direcciones previstas antes de autorizar el
\code{apply}.

\begin{figure}[H]
  \centering
  \includegraphics[width=\textwidth]{01b-terraform-plan-resumen.png}
  \caption{Resumen del plan Terraform de Fase 2. La salida acredita 41 altas sin cambios ni
  destrucciones sobre el estado utilizado en la ventana.}
  \label{fig:f2-terraform-plan}
\end{figure}

El \code{apply} produjo identificadores reales para la VPC y las subredes, las IP privadas de los
clientes y una IP pública temporal para Visitas. La captura siguiente documenta los outputs usados
por la suite; los archivos de estado y el plan binario no se versionaron.

\begin{figure}[H]
  \centering
  \includegraphics[width=0.98\textwidth]{02-terraform-apply-outputs.png}
  \caption{Outputs obtenidos después del despliegue. Se observan las subredes, los clientes de TI
  y Visitas y las direcciones utilizadas para las pruebas.}
  \label{fig:f2-terraform-apply}
\end{figure}

\begin{infobox}[Trazabilidad]
Las capturas acreditan el plan y los outputs del despliegue; los resultados funcionales se
demuestran por separado con \code{evidence/fase-2/pruebas-fase2.csv} y los archivos individuales
F2-01 a F2-13. Esta separación evita presentar el \code{apply} como sustituto de las pruebas.
\end{infobox}
